% =============================================================================
% Section: Runtime Overhead Characterization (Experiment 1)
% Paste into the Evaluation section of the SC paper.
%
% Required figures (drop into figs/exp1/):
%   exp1_overhead_stacked.pdf
%   exp1_crossover.pdf
%
% Cross-references assumed to exist elsewhere:
%   \label{sec:observability}    - Section 3.4
%
% Outstanding TODO: replace \TODO{platform} with the testbed name.
% =============================================================================

\subsection{Runtime Overhead Characterization}
\label{sec:exp-runtime-overhead}

\paragraph{Question.}
When the user's callback compute time changes, which \drava{} runtime
components account for the remaining wall time, and how do their costs
scale?  We use this experiment to attribute end-to-end latency to specific
runtime mechanisms rather than to ``the framework'' as a whole.

\paragraph{Method.}
We run the two-stage PtychoNN pipeline with both per-stage callback batch
thresholds matched to the inference batch size, $B^{cb}_1 = B^{cb}_2 = B$,
sweeping $B \in \{32, 64, 128, 256, 512\}$.  Varying $B$ changes the
amount of user work per callback (a larger inference batch is a larger
GPU call) and simultaneously changes the size of each unit of work the
runtime handles, so the experiment exercises every runtime component
under five different operating points.  Per-stage worker counts are
fixed at $N_1 = N_2 = 4$, the fetch timeout at $\tau = 200$\,ms, and the
publisher rate is uncapped ($r = 0$) so that the runtime is the sole
pacer.  Each configuration processes $N = 10{,}000$ frames on
\TODO{platform}.

For each run we use the per-vertex counters from
Section~\ref{sec:observability} to split each stage's wall time into
four non-callback components: \emph{transport} (JetStream publish wait on
the egress side), \emph{dispatch} (per-batch handoff into the XKRT
work-stealing queue), \emph{microbatch wait} (time spent in
timeout-driven flush), and \emph{publish} (callback-internal publish
calls).  We deliberately exclude the callback compute time from the
comparison because it is reported by Stage~1 as cumulative across the
four worker threads and by Stage~2 as per-batch wall time, and the two
quantities are not directly comparable.

\paragraph{Results.}
Table~\ref{tab:exp1-summary} reports the end-to-end latency $T_{e2e}$ and
the two largest non-callback components for each batch size.  $T_{e2e}$
is non-monotonic in $B$ ($9.97 \to 6.42 \to 4.09 \to 7.09 \to 13.99$\,s);
the minimum at $B{=}128$ is $2.44\times$ faster than $B{=}32$ and
$3.42\times$ faster than $B{=}512$.  The publisher throughput stays flat
at $\approx 17$\,kfps across the sweep, so the pipeline is bounded
entirely by the two stages.

\begin{table}[t]
\centering
\caption{End-to-end pipeline and dominant per-stage runtime costs at
$N_1{=}N_2{=}4$, $\tau{=}200$\,ms, $r{=}0$, $N{=}10{,}000$\,frames.
Throughput in frames/s; times in s.  $T_{\text{S}_1}^{\text{tr}}$ is
Stage~1 transport cost; $T_{\text{S}_2}^{\text{dsp}}$ is Stage~2 dispatch
cost.}
\label{tab:exp1-summary}
\small
\begin{tabular}{rrrrrr}
\toprule
$B$ & $\text{S}_1$ fps & $\text{S}_2$ fps & $T_{\text{S}_1}^{\text{tr}}$ & $T_{\text{S}_2}^{\text{dsp}}$ & $T_{e2e}$ \\
\midrule
 32 & 1{,}054 & 1{,}095 &  0.48 &  8.61 &  9.97 \\
 64 & 1{,}685 & 1{,}803 &  0.49 &  4.49 &  6.42 \\
128 & 2{,}801 & 3{,}213 &  0.52 &  1.31 &  4.09 \\
256 & 4{,}216 & 1{,}761 &  4.72 &  0.12 &  7.09 \\
512 & 5{,}506 &   840   & 12.17 &  0.03 & 13.99 \\
\bottomrule
\end{tabular}
\end{table}

Figure~\ref{fig:exp1-breakdown} shows the four-component decomposition as
stacked bars, one panel per stage.  The decomposition is concrete: two
components dominate, and the remaining two are negligible at every
sampled $B$.

\begin{itemize}
  \item Stage~1's non-callback cost is essentially all \emph{transport}.
        It grows from $0.48$\,s at $B{=}32$ to $12.17$\,s at $B{=}512$,
        a $25\times$ increase, because each callback emits a payload
        proportional to $B$ and the JetStream publish wait scales with
        payload size.
  \item Stage~2's non-callback cost is essentially all \emph{dispatch}.
        It falls from $8.61$\,s at $B{=}32$ to $0.03$\,s at $B{=}512$,
        a $287\times$ decrease, because the fixed per-batch overhead of
        NumPy reconstruction and XKRT submission amortises across
        larger batches.
  \item \emph{Microbatch wait} stays below $10^{-6}$\,s for every
        $(B, \text{stage})$ cell, confirming that the callback-size
        flush condition fires before the timeout in this regime.
  \item \emph{Publish} stays below $0.1$\,s on Stage~1 and is exactly
        zero on Stage~2 (the stitching callback does not republish).
\end{itemize}

Figure~\ref{fig:exp1-crossover} plots the two dominant components on
shared log-log axes.  They are monotonic in opposite directions and
cross at $B \approx 142$, which is the same operating point at which
$T_{e2e}$ reaches its empirical minimum at $B{=}128$ (the nearest
sampled value).  The two runtime mechanisms therefore explain the
U-shape in $T_{e2e}$ directly: at $B \leq 64$ the pipeline is
dispatch-bound on Stage~2, at $B \geq 256$ it is transport-bound on
Stage~1, and the optimum is the batch size at which the two costs are
balanced.

\paragraph{Takeaway.}
The runtime overhead at the optimal $B{=}128$ is dominated by a single
mechanism on each stage (Stage~1: $0.52$\,s of transport; Stage~2:
$1.31$\,s of dispatch); together they account for $1.83$\,s of the
$4.09$\,s pipeline latency, or $44.8\%$.  Reducing $T_{e2e}$ below the
$4.09$\,s floor measured here therefore requires targeting one of these
two specific components --- shrinking the egress payload (e.g. by
quantising Stage~1 outputs) or eliminating the per-batch dispatch cost
on Stage~2 (e.g. by fusing the stitching kernel with the work-stealing
runtime) --- rather than tuning configuration knobs in isolation.

\begin{figure}[t]
  \centering
  \includegraphics[width=\columnwidth]{figs/exp1/exp1_overhead_stacked.pdf}
  \caption{Decomposition of per-stage non-callback runtime cost into the
  four components emitted by the runtime's per-vertex counters, as a
  function of $B$.  Stage~1 cost is essentially all transport (blue);
  Stage~2 cost is essentially all dispatch (amber); publish (green) and
  microbatch wait (purple) are negligible at every $B$.}
  \label{fig:exp1-breakdown}
\end{figure}

\begin{figure}[t]
  \centering
  \includegraphics[width=\columnwidth]{figs/exp1/exp1_crossover.pdf}
  \caption{The two dominant components from Figure~\ref{fig:exp1-breakdown}
  on log-log axes.  Stage~1 transport cost rises with $B$ while Stage~2
  dispatch cost falls; they intersect at $B \approx 142$, which
  coincides with the empirical minimum of $T_{e2e}$ at $B{=}128$.}
  \label{fig:exp1-crossover}
\end{figure}
